% Optional introduction/footnote material, status checked 8 October 2026.
% Bibliography keys are supplied in su2_sources.bib.

\paragraph{Current literature and the meaning of the target.}
The present work does not claim a new general quasipolynomial algorithm.
Lackenby's March 2021 author slides announce a running time
$2^{O((\log n)^3)}=n^{O((\log n)^2)}$ and outline a hierarchy approach
\cite{Lackenby2021Slides}. We use that source's displayed bound.
His July 2026 preprint proves a new hierarchy
algorithm for incompressibility, with unknot recognition as a consequence,
and explicitly leaves possible speed-ups to future work
\cite{Lackenby2026Hierarchies}.  It should therefore not be cited as a
completed proof of the announced quasipolynomial bound.  The present
algebraic sufficient condition targets $n^{O(\log n)}$ on a specified
residual class, under explicitly stated and verified compression hypotheses.
These are different statements.

A September 2026 preprint by Musick claims a polynomial-time procedure
for producing locally minimal bridge presentations and, consequently,
unknot recognition in $P$ \cite{Musick2026}.  Version 2 is dated
20 September 2026.  That claim has not been independently established by
this study and none of the results or code here depends on it.  Recording
its existence avoids the inaccurate claim that no polynomial-time preprint
has appeared; treating it as an unverified recent claim avoids replacing
proof checking with an abstract-level assertion.  An independent audit
would have to establish the completeness of the permitted presentation
moves and the polynomial bounds on their compressed execution, as well as
the correctness of the initial diagram conversion.

\paragraph{What the hierarchy bounds suggest.}
Lackenby's 2026 Proposition 9.1 bounds the number of main iterations by
$L(g+1)^L$, where $L$ controls hierarchy length and $g$ controls the
relevant pattern complexity.  The later sections develop efficient
hierarchy certificates and compressed cutting operations
\cite{Lackenby2026Hierarchies}.  Purely as an implication of that bound,
polynomial $g$ and $L=O(\log n)$ would yield $n^{O(\log n)}$
iterations, while $L=O((\log n)^2)$ yields
$n^{O((\log n)^2)}$.  Counting iterations alone still requires a bound
on the bit cost of each iteration and on the construction of the
surfaces used.  The 2021 slides describe polynomial genus, compressed
hierarchies, multisurfaces, and Heegaard/Cheeger-region arguments as the
components of that route \cite{Lackenby2021Slides}.  A rigorous
implementation project would need exact data formats, verified transition
rules, and complete size and bit-complexity bounds for those components.
